% Transcribed from sections/*.md. The Markdown is the source of truth:
% edit there, then run build.sh, which regenerates this file.
\documentclass[10pt,twocolumn]{article}
\usepackage[T1]{fontenc}
\usepackage{lmodern}
\usepackage[margin=0.75in]{geometry}
\usepackage{xcolor}
\usepackage{hyperref}
\newcommand{\TBD}[1]{\textcolor{red}{[TBD: #1]}}
\newcommand{\pilot}{\textcolor{orange}{[pilot]}}
\newcommand{\chk}[1]{\textcolor{blue}{[check: #1]}}

\title{Images Are Unique, Blocks Are Not:\\Sharing Disk State Across a Fleet of Agent VMs}
\author{Anonymous}
\date{}

\begin{document}
\maketitle
\begin{abstract}
AI agents, reinforcement learning environments (RL) and continuous integration pipelines (CI) often run independent tasks in isolated virtual machines (VMs), booting their own disk images. Published deployments keep tens of thousands of such images, terabytes to over a hundred terabytes in total, many built on a handful of base images. Operators reuse disk data by image identity, not by content. A host skips fetching only when it already holds that exact image. Some place a VM by image locality, while others place by load and accept the cost of fetching. We ask how much content distinct images share at the block level. We captured the disks of [N] VMs from RL environments, coding-agent sessions and CI jobs, comparing them block by block. [X\%] of the bytes a task writes, and [R\%] of the blocks it reads before its first command, are identical to blocks on the disk of a VM it was not copied from. A host can use those blocks only if it holds them, which is the case for [Y\textsubscript{1}\%] to [Y\textsubscript{2}\%] depending on placement. A content-addressed store names each block by its fingerprint, so a host recognises the blocks it already holds under another image and fetches only the rest. We measure such a store against ZFS deduplication on each host, on the same machines. [Result.] We release the corpus.
\end{abstract}

%% ---- sections/01-introduction.md ----
\section{Introduction}

Each task an AI agent runs gets its own environment image. In 2024, SWE-Gym shipped one image per task, 2,438 images and 6 TB in total \cite{swe-gym}. Moonshot reports creating tens of thousands of VMs within seconds, each with a unique set of images \cite{kimi-k3}. In one week of early 2026, DeepSeek's platform held more than 130 TB of environment images, and its microVMs alone used two base images and 53,590 task-specific disks \cite{dsec}. Moving these images is costly. Each of DeepSeek's VMs reads only 4.2\% to 13.3\% of its image, and loading whole images before start made a burst of VMs take 1.71 times as long \cite{dsec}. Alibaba reports that when environments time out, resetting them takes 78\% of rollout time, with image downloads saturating the network \cite{rollart}.

Operators respond in two ways. Some decide where a VM runs by what the host already holds. Cloudflare favors hosts that have the VM's image or snapshot \cite{cloudflare-sandboxes}. Prime Intellect schedules VMs where their environment is already present \cite{prime-sandboxes}. Kubernetes scores nodes by the images they hold \cite{k8s-imagelocality}. Others place by load and fetch what is missing. DeepSeek samples a few hosts and picks the least loaded \cite{dsec}. Sandman spreads VMs across hosts, and reports that concentrating them on hosts with the image cached caused failures to accumulate \cite{sandman}.

Both responses treat the image as the unit. Placing by image helps only when the same image is used again, and DeepSeek measures that a microVM image is typically used by one VM within a job \cite{dsec}. Placing by load gives up on reuse at the host.

Sharing below the image has been reported only for files. DeepSeek builds container environments from shared layers and serves the files accessed by 70\% of its container tasks from one 30 TB deduplicated set \cite{dsec}. A VM's disk is different. It is a block device holding a file system, written to as the task runs, and DeepSeek stores each microVM workspace as its own disk image \cite{dsec}. For such disks we found no report of how many blocks two different VMs have in common.

This paper measures that. A \emph{task} is one unit of agent work run in one VM, which may be an RL episode, a coding-agent session or a CI job. We ask three questions.

\begin{enumerate}
\item \textbf{Block fanout.} How much of what a task writes, and of what it reads before its first command, is identical to blocks on the disk of a VM it was not copied from? (\S4)

\item \textbf{Placement.} How many of those duplicate blocks are already on the host, under placement by load, by image and by repository? (\S5)

\item \textbf{Cost.} What does it cost to share them, with deduplication on each host or with one content-addressed store across all hosts? (\S7, \S8)

\end{enumerate}

A content-addressed store names each block by its fingerprint. A host can then recognise blocks it already holds under another image and fetch only the rest. Fetching on demand is established practice \cite{lambda,dsec}. What we measure is how much the recognition adds to it.

We answer the questions on the captured disks of \TBD{N} VMs from three workloads and on \TBD{2 or 4} bare-metal hosts.

\textbf{Findings.} \TBD{one sentence per question, written when the data exists.}

\textbf{Contributions.}

\begin{itemize}
\item A released corpus of VM disks from three agent workloads, captured at block level with file attribution.

\item A measurement of block fanout across VMs that share no snapshot.

\item The split of duplicate blocks between the same host and other hosts under each placement policy, with the load imbalance each causes.

\item A matched comparison of per-host ZFS deduplication and a fleet-wide content-addressed store.

\end{itemize}

The content-addressed store is the system under test. Its mechanism is not new (\S9) and we do not claim it.

%% ---- sections/02-background.md ----
\section{Background}

\textbf{Provisioning a VM disk.} A VM needs its image on the host before it can start. The host either copies the whole image or loads it on demand, fetching each part when it is first read \cite{slacker,lambda}. A \emph{cold start} is a start on a host that holds none of the image.

\textbf{Sharing by lineage.} A \emph{snapshot} is a read-only copy of an image at one moment. A \emph{clone} is a writable image whose initial contents equal a snapshot \cite{openzfs-concepts}. Two disks \emph{share lineage} if one is a clone of a snapshot of the other, or both are clones of a common snapshot. Snapshots, clones and image layers share only content that lineage records. Two VMs that install the same package independently share none.

\textbf{Sharing by content.} \emph{Deduplication} stores one copy of equal data. The store splits data into \emph{chunks}, computes a \emph{fingerprint} (a collision-resistant hash) of each, and keeps one copy per fingerprint \cite{venti,lbfs}. A store whose chunks are addressed by fingerprint is \emph{content-addressed}. Deduplication finds equal content regardless of lineage.

\textbf{Deduplication on one host.} ZFS deduplicates blocks inline, on the write path, against a deduplication table \cite{openzfs-concepts}. The table belongs to one storage pool, and a pool belongs to one host \chk{cite}. Equal blocks on two hosts are stored twice, and moving a disk between hosts sends them again \chk{\texttt{zfs send} behaviour}.

\textbf{Deduplication across a fleet.} A fleet-wide content-addressed store keeps one fingerprint namespace for all hosts. AWS Lambda stores container images this way and loads chunks on demand \cite{lambda}. Such a store pays for an index, for a second write when deduplication is post-process, and for a network round trip when a host reads a chunk it does not hold.

\textbf{Placement decides what a host can share.} With per-host deduplication, two equal chunks are shared only if their VMs run on the same host. Placing similar virtual machines together to raise host-local sharing is an old idea, studied for memory pages \cite{memory-buddies,sindelar} and for disk blocks \cite{smio}. It trades sharing against load balance \cite{dong}. VM operators today place by load (\S1), which is the opposite end of that trade.

%% ---- sections/03-workloads.md ----
\section{Workloads and Corpus}

We study three workloads, chosen because each creates many VMs from related but distinct images.

\textbf{RL rollout environments.} One VM per episode, started from a per-task image and discarded. We use SWE-bench Verified (500 tasks, 12 repositories) \cite{swebench-verified} and Terminal-Bench 2 \TBD{count} \chk{cite}. Creation arrives in bursts of many distinct images.

\textbf{Coding-agent sessions.} The same disk contents as an RL episode, but long-lived: the session idles, is suspended, and resumes, possibly on another host. We model this with the same tasks, paused after setup. \TBD{confirm workload definition}

\textbf{CI jobs.} One VM per job, built from a runner image and a repository checkout. \TBD{N} jobs from \TBD{N} repositories across \TBD{languages}. We expect less duplication here and include it as a contrast.

\textbf{Capture.} Each task runs in a VM with a 16 GiB ext4 disk. We capture the whole disk at fixed points: after setup and after the task for coding tasks, after the job for CI. Free blocks are zeroed using the file system's allocation map so that deleted data is not counted. \chk{against the capture code}

\textbf{Counting.} We fingerprint each captured disk at 4 KiB and 16 KiB fixed-size chunks and with content-defined chunking. We count only non-zero chunks that are absent from the VM's base image, so that sharing by lineage is excluded by construction. A file manifest (path, size, extents) attributes each chunk to a category: dependency cache, installed dependencies, toolchain, build output, repository, system.

\textbf{Corpus.} \TBD{table. Workload, tasks, disks, bytes written beyond the base. Pilot: 10 coding tasks, 15 CI jobs.} \pilot{} We release fingerprints, manifests and capture scripts. \TBD{what cannot be released}

%% ---- sections/04-block-fanout.md ----
\section{Block Fanout}

The \emph{block fanout} of a chunk is the number of VMs whose disks contain it. A chunk a task writes is a \emph{duplicate} if an earlier VM, sharing no lineage with it, already contains the same chunk; otherwise it is \emph{unique}.

\textbf{How much is duplicated.} \TBD{fraction of written bytes that are duplicates, per workload, as a distribution over tasks. Pilot: coding tasks 10.2 of 15.1 GB (68\%); CI 3.0 of 12.8 GB (24\%).} \pilot{}

\textbf{Fanout distribution.} \TBD{CDF of block fanout per workload, beside the published image fanout.}

\textbf{What the duplicates are.} \TBD{duplicate bytes by path category. Pilot: for coding tasks, 9.5 of 12.4 GB written to dependency caches duplicate an earlier task; for CI, build output (4.7 GB) is almost all unique.} \pilot{}

\textbf{Within and across repositories.} \TBD{Pilot: 12.1 of 13.2 GB of duplicates are within one repository.} \pilot{}

\textbf{Chunking.} \TBD{fixed 4 KiB, fixed 16 KiB and CDC. Pilot: 13.2, 11.8 and 14.3 GB duplicated.} \pilot{}

\textbf{Simpler alternatives.} Much duplicate content is package data, which a shared package cache or one image per repository would also share. \TBD{bytes each alternative leaves duplicated, on the same corpus.}

\textbf{Summary.} \TBD{one sentence.}

%% ---- sections/05-placement.md ----
\section{Where Duplicates Sit}

A duplicate chunk helps a host only if the host can reach the earlier copy. A \emph{same-host duplicate} has an earlier copy on the same host. An \emph{other-host duplicate} has copies only on other hosts. Per-host deduplication shares the first kind; only a fleet-wide store shares the second.

\textbf{Method.} We assign the captured VMs, in arrival order, to H hosts under a placement policy, and classify each duplicate chunk. Policies:

\begin{itemize}
\item \emph{load-based}: the least-loaded host, as operators report doing \cite{dsec};

\item \emph{random}: a uniformly random host;

\item \emph{affinity}: a host that already ran a VM of the same repository;

\item \TBD{one or two points between load-based and affinity}.

\end{itemize}

For each policy we report the same-host and other-host shares of duplicate bytes and the load imbalance across hosts. \TBD{imbalance metric}

Placement at 4, 16 and 64 hosts is simulated over captured disks. We check the simulation against a real run on four hosts. \TBD{only if the four-host reservation is approved}

\textbf{Results.} \TBD{figure, same-host share against H per policy. Pilot, 4 KiB: affinity keeps 12.1 of 13.2 GB same-host at H = 4 and 16; random keeps 8.8 GB at H = 4 and 2.3 GB at H = 16.} \pilot{}

\textbf{Load imbalance.} \TBD{In SWE-bench Verified one repository holds 231 of 500 tasks, so strict affinity concentrates load.}

\textbf{Summary.} \TBD{one sentence.}

%% ---- sections/06-system.md ----
\section{The System}

We measure a content-addressed block backend for VM disks. It is the system under test; this section describes only what the measurements depend on.

\textbf{Interface.} One host process serves private disks to unmodified QEMU guests as virtio-blk devices over vhost-user. Content addressing is invisible to the guest.

\textbf{Write path.} Each image has a write-ahead log (WAL). A write is appended to the log; a guest FLUSH is acknowledged after \texttt{fdatasync} on the local device. Deduplication is post-process: a background compactor reads logged writes, fingerprints each 4 KiB chunk with BLAKE3, stores chunks not already present in a chunk store shared by all images on the host, and publishes the image's new manifest. Logged data is reclaimed after it is compacted.

\textbf{Read path.} A read is served from the log if the block was written since the last compaction, and otherwise from the manifest and the chunk store.

\textbf{Images.} A \emph{manifest} maps each offset of an image to a chunk fingerprint. A snapshot is a retained manifest; a clone is a copy of one.

\textbf{Across hosts.} Hosts form a static fleet over kernel TCP. Each chunk is stored on k hosts chosen by its fingerprint (k = 1 or 2). A host that reads a chunk it does not hold fetches it by fingerprint from a host that stores it and verifies the fingerprint before use. Provisioning or moving an image sends its manifest; chunks follow on demand.

\textbf{Limits.} FLUSH durability is local: data acknowledged but not yet replicated is lost with its host. Membership is static. Memory state is not handled. \TBD{implementation size, language, revision.}

%% ---- sections/07-cost.md ----
\section{What Sharing Costs}

\textbf{Testbed.} \TBD{hosts, CPU, memory, NVMe, network, kernel, guest shape. Testbed floor: raw 4 KiB read, \texttt{fdatasync}, TCP round trip.}

\subsection{One host}

We run the same workload on three backends on the same NVMe device: the content-addressed store, ZFS with deduplication at 4 KiB records, and a raw file on XFS with no sharing.

\begin{itemize}
\item \textbf{Bytes stored.} \TBD{}

\item \textbf{Latency.} \TBD{read, write and FLUSH, p50 and p99, at 1, 4, 8 and 16 guests.}

\item \textbf{Memory.} \TBD{fingerprint index against the ZFS deduplication table.}

\item \textbf{Write amplification.} \TBD{}

\end{itemize}

\subsection{Across hosts}

\begin{itemize}
\item \textbf{A remote read.} \TBD{guest-visible latency of a read whose chunk is on another host, from that host's memory and from its disk, against a local read.}

\item \textbf{Cold-start burst.} \TBD{N distinct images started at once on a host that holds none of them: time to first command, bytes fetched, fraction of the image read.}

\item \textbf{Resuming on another host.} \TBD{time and bytes to resume a paused VM elsewhere, against a full copy and an incremental \texttt{zfs send}. Where the two disks share lineage, incremental send moves fewer bytes; we report it.}

\item \textbf{Hedged requests.} \TBD{p99 with and without a hedged request, and the extra bytes sent. Keep only if the effect is clear.}

\end{itemize}

\textbf{Summary.} \TBD{one sentence per subsection.}

%% ---- sections/08-crossover.md ----
\section{Per-Host Deduplication or a Fleet Store}

Sections 5 and 7 combine into one comparison. For each placement policy we compute, for per-host ZFS deduplication and for the fleet store:

\begin{itemize}
\item bytes stored across the fleet;

\item bytes moved to start each VM;

\item time to first command;

\item memory spent on indexes;

\item load imbalance.

\end{itemize}

\TBD{figure. One panel per metric, placement policy on the x-axis, one line per design.}

\textbf{Where they cross.} \TBD{the placement policy and fleet size beyond which the fleet store stores or moves fewer bytes, and what it pays in latency there.}

\textbf{What an operator should do.} \TBD{stated as conditions, not as a recommendation for one design.}

%% ---- sections/09-related-work.md ----
\section{Related Work}

\textbf{Content-addressed storage and deduplication.} Venti addresses blocks by fingerprint \cite{venti}; LBFS adds content-defined chunking \cite{lbfs}. Deduplication of primary data and of VM images has been measured at scale \cite{meyer,elshimi,jin-miller,jayaram}. Those studies analyse stored images, most sharing an operating-system lineage. We measure blocks written by independent tasks.

\textbf{Moving VM state on demand.} Fetching a VM's disk on demand at another host, and skipping data the destination already holds, dates to \cite{sapuntzakis,isr,snowflock}. AWS Lambda loads container images on demand from a fleet-wide content-addressed store \cite{lambda}. We reuse these mechanisms and claim none.

\textbf{Placement and sharing.} Co-locating VMs to raise host-local sharing was studied for memory \cite{memory-buddies,sindelar} and for disk blocks \cite{smio}. Routing data among deduplicating nodes trades deduplication against skew \cite{dong}. We apply the same trade to VM disks and add the fleet-store alternative on the same axes.

\textbf{VM infrastructure.} DSec reports image fanout and places by load \cite{dsec}; Sandman rejects placement by image \cite{sandman}. Reports on SWE-bench image sizes work at layer or file level \cite{epoch-swebench,logicstar}. We add the block level and the placement analysis.

\textbf{Comparing deduplicating stores.} DupHunter compares ZFS and Ceph deduplication on Docker registry layers, with unmatched block sizes \cite{duphunter}. We compare on VM disks at equal chunk size.

\chk{Sun et al. (ACM TOS 2018) and Graviet et al. (2026) are unread and may narrow the placement claim.}

%% ---- sections/10-limits.md ----
\section{Limits}

\begin{itemize}
\item \textbf{Fleet size.} Latency is measured on \TBD{2 or 4} hosts. Placement at larger fleets is simulated over captured disks.

\item \textbf{Disk state only.} We do not capture or restore memory. Operators report that memory, not disk, bounds VM density.

\item \textbf{Workloads.} RL environments are Python repositories from two benchmarks. Other languages and proprietary workloads may share less.

\item \textbf{Arrival order.} A duplicate is defined against earlier VMs, so the result depends on the order tasks arrive. \TBD{sensitivity}

\item \textbf{Network.} Kernel TCP only. RDMA would lower the cost of a remote read.

\item \textbf{Durability.} The fleet store acknowledges FLUSH locally; we do not evaluate replicated durability.

\item \textbf{Startup in context.} For long tasks, seconds of startup matter little. The results bear on short tasks and bursts.

\end{itemize}

\bibliographystyle{plain}
\bibliography{refs}
\end{document}
